Este trabalho apresenta o projeto de desenvolvimento de uma arquitetura denominada "EDGAr", que integra um motor de Geração Aumentada por Recuperação Multimodal (MRAG) destinado a otimizar a compreensão de documentos no cenário educacional. O sistema é composto por pipelines de indexação híbrida e algoritmos de extração que realizam buscas de granularidade fina para verificar e isolar componentes textuais e visuais estruturados antes de formular respostas. O fluxo de dados entre os módulos ocorre via representações vetoriais, e as informações recuperadas alimentam a aplicação cliente "LuminIA Reader", que atua como camada de aplicação. O projeto visa prevenir a ocorrência de alucinações causadas por ruídos em documentos longos, garantindo contextos densos com baixo custo computacional. Os testes de arquitetura planejados buscam demonstrar a leveza do pipeline, validando sua funcionalidade e viabilidade para futuras implementações em um protótipo funcional. Esta arquitetura propõe-se como uma contribuição relevante para a tecnologia educacional, viabilizando interações mais precisas entre o estudante e o material didático ao mitigar o problema das alucinações semânticas.

% Separe as palavras-chave por ponto
\palavraschave{Geração Aumentada por Recuperação; Visão Computacional; Materiais Didáticos; Granularidade Fina.}